



\textbf{Problem 11.2}

\begin{enumerate}[label={(\alph*)}]
\item By definition we have that
\[
	\int_B \bbE[X | \cH] \, \dd \bbP = \int_B X \, \dd \bbP,
\]
holds for all $B \in \cH$. Since by assumption both $\bbE[X | \cH]$ and $X$ are $\cH$-measurable, the result follows from problem 8.2.
\item Note that $a \bbE[X | \cH]$ is $\cH$-measurable. Moreover,
\[
	\int_B a \bbE[X | \cH] \, \dd \bbP = a \int_B \bbE[X | \cH] \, \dd \bbP = a \int_B X \, \dd \bbP = \int_B aX \, \dd \bbP.
\]
This proves the claim.
\item Similarly to the previous point, we first note that since $\bbE[X | \cH]$ and $\bbE[Y | \cH]$ are $\cH$-measurable so is $\bbE[X | \cH] + \bbE[Y | \cH]$. The result then follows because
\begin{align*}
	\int_B \bbE[X | \cH] + \bbE[Y | \cH] \, \dd \bbP 
	&= \int_B \bbE[X | \cH] \, \dd \bbP + \int_B \bbE[Y | \cH] \, \dd \bbP \\
	&= \int_B X \, \dd \bbP + \int_B Y \, \dd \bbP = \int_B X + Y \, \dd \bbP.
\end{align*}
\item First we observe that for any $B \in \cH$
\[
	\int_B \bbE[X | \cH] \, \dd \bbP = \int_B X \, \dd \bbP \le \int_B Y \, \dd \bbP = \int_B \bbE[Y | \cH] \, \dd \bbP.
\]
Now consider the event $A := \{ \bbE[X | \cH] > \bbE[Y | \cH]\} \in \cH$. If this event has non-zero measure then it would follow that
\[
	\int_A \bbE[X | \cH] \, \dd \bbP > \int_A \bbE[Y | \cH] \, \dd \bbP,
\]
which is a contradiction. Hence we conclude that $\bbE[X | \cH] \le \bbE[Y | \cH]$ holds $\bbP$-almost everywhere.
\end{enumerate}

\bigskip

\textbf{Problem 11.3}

\begin{enumerate}[label={(\alph*)}]
\item This follows by repeating the step for the solution to Problem 4.8 a). 
%
%By definition, we have that $\nu_X(\Omega) = \int_\Omega X\,\dd\mu = 1$. Now let $(A_n)_{n\in\bbN}$ be a family of mutually disjoint measurable sets in $\sigma(X)$. Then we have that the sequence 
%	\[
%		g_n:= \sum_{i=1}^n X\,\mathbf{1}_{A_i} = X\,\mathbf{1}_{\bigcup_{i=1}^n\! A_i}\;\longrightarrow\; g:= X\,\mathbf{1}_{\bigcup_{i\in\bbN}\! A_i}\qquad\text{pointwise monotonically}.
%	\]
%	By MCT, we then have that
%	\[
%		\nu_X\left(\bigcup_{i\in\bbN} A_i\right) = \int_{\bigcup_{i\in\bbN}\! A_i} X\,\dd\mu = \lim_{n\to\infty} \int_{\bigcup_{i=1}^n\! A_i} X\,\dd\mu = \lim_{n\to\infty} \sum_{i=1}^n \int_{A_i} X\,\dd\mu = \sum_{i\in\bbN}\nu_X(A_i),
%	\]
%	thus showing that $\nu_X$ is a measure on $(\Omega,\sigma(X))$.

\item We start by observing that the result will directly follow from Theorem 2.15 if we can show that the \sigalg/ $\sigma(X)$ satisfies the two properties.

The first one is immediate, from the fact that $X^{-1}(A) \cap X^{-1}(B) = X^{-1}(A\ cap B)$. For the second one consider the intervals $I_n = (-n,n)$ and define $A_n := X^{-1}(I_n) \in \sigma(X)$. Since $\bigcup_{n \in \bbN} I_n = \bbR$ it follows that $\bigcup_{n \in \bbN} A_n = \Omega$. Moreover since $|X| \le n$ on the set $A_n$ it holds that
\[
	\mu(A_n) = \nu_X(A_n) = \int_{A_n} X \, \dd \bbP \le n \int_{A_n} \, \dd \bbP = n \bbP(A_n) \le n < \infty.
\]
Thus the second condition of Theorem 2.15 is also satisfied and the result now follows.

	
\end{enumerate}

\bigskip

\textbf{Problem 11.4}

The solution to this problem will closely follow the proof of Lemma 11.7. To this end let $g(x,y)$ denote the conditional density of $X$ given $Y = y$ and define the function $\phi(y) := \int_\bbR h(x) g(x,y) \, \lambda(\dd x)$. We now need to show that for any $A \in \cB_\bbR$ it holds that
\[
	\int_{Y^{-1}(A)} \phi(Y) \, \dd \bbP = \bbE[\mathbbm{1}_A(Y) h(X)].
\]

Using the change of variables formula and the definition of $\phi$ and $g$ we get
\begin{align*}
	\int_{Y^{-1}(A)} \phi(Y) \, \dd \bbP &= \int_\Omega \mathbbm{1}_{Y^{-1}(A)} \phi(Y) \, \dd \bbP\\
	&= \int_\bbR \int_A(y) \phi(y) f_Y(y) \, \lambda(\dd y)\\
	&= \int_{\bbR} \int_\bbR \mathbbm{1}_A(y) h(x) g(x,y) f_Y(y)  \, \lambda(\dd x) \, \lambda(\dd y)\\
	&= \int_{\bbR^2} \mathbbm{1}_{A}(y) h(x) f(x,y) \, \lambda(\dd x) \, \lambda(\dd y)\\
	&= \bbE[\mathbbm{1}_{A}(Y) h(X)].
\end{align*}



\textbf{Problem 11.5}
\begin{enumerate}[label={(\alph*)}]
\item First note that for any $\ell \ge 0$ 
\[
	\bbP(X > \ell) = \int_{\ell} \nu e^{-\nu x} \, \lambda(dx) 
	= e^{-\nu \ell}.
\]

By Lemma 11.5 we have that
\[
	\bbP(X > t +s | X > t) = \frac{\bbP(X > t + s, X > t)}{\bbP(X > t)}
	= \frac{\bbP(X > t + s)}{\bbP(X > t)}
	= \frac{e^{-\nu(t + s)}}{e^{-\nu t}} = e^{-\nu s} = \bbP(X > s).
\]
\item Define the function $f(x,y) = \mathbbm{1}_{x + y \le t}$. By Lemma 11.8 we then have that
\[
	\bbP(X + Y \le t | Y = y) = \bbE[f(X,Y) | Y = y]
	= \bbE[f(X,y)].
\]
Using the probability density function we then get
\[
	\bbE[f(X,y)] = \int_{\bbR} f(x,y) \rho(x) \, \lambda(dx)
	= \mathbbm{1}_{y < t} \int_0^{t-y} \nu e^{-\nu x} \, \lambda(dx)
	= \mathbbm{1}_{y < t} (1 - e^{-\nu(t-y)}).
\]
\item We have that
\begin{align*}
	\bbP(X + Y \le t) &= \bbE[\bbP(X+Y \le t | Y)]\\
	&= \int_\bbR \rho(y) \bbP(X + Y \le t | Y = y) \, \lambda(dx)\\
	&= \int_\bbR \mathbbm{1}_{0 \le y < t} (1 - e^{-\nu(t-y)}) \nu e^{-\nu y} \, \lambda(dy) \\
	&= \int_0^t \nu e^{-\nu y} - e^{-\nu t} \, \lambda(dy) \\
	&= 1 - e^{-\nu t} - t e^{-\nu t}.
\end{align*}
\end{enumerate}


\bigskip

\textbf{Problem 11.6}

\begin{enumerate}[label={(\alph*)}]
%\item First note that it always holds that
%\[
%	\int_\Omega Z \, \dd \bbP = \int_{\Omega \times \Omega} f \circ (X,Y) \, \dd \Delta_\# \bbP.
%\]
%
%Furthermore, as shown in the proof of Lemma 11.8 if holds that
%\[
%	\Delta_\# \bbP = \bbP \otimes \bbP
%\]
%on the set $X^{-1}(\cB_\bbR) \times Y^{-1}(\cB_\bbR)$. Note that by Lemma 5.2 we have that
%\[
%	\sigma(X) \otimes \sigma(Y) = \sigma(X^{-1}(\cB_\bbR) \times Y^{-1}(\cB_\bbR)). 
%\]
%Now let $I_n = (-n,n)$. Then since both $\Delta_\# \bbP$ and $\bbP \otimes \bbP$ are finite measures, and the sequence $A_n := X^{-1}(I_n) \times Y^{-1}(I_n)$ satisfies $\bigcup_{n \ge 1} A_n = \Omega \times \Omega$, it follows from Theorem 2.15 that $\Delta_\# = \bbP \otimes \bbP$ on $\sigma(X) \otimes \sigma(Y)$.
%
%Suppose now that $f$ is a simple function, i.e. $f = \sum_{i = 1}^N f_i \mathbbm{1}_{A_i}$ with $A_i \in \cB_{\bbR^2}$. Then $f \circ (X,Y) = \sum_{i = 1}^N f_i \mathbbm{1}_{(X,Y)^{-1}(A_i)}$ and since $(X,Y)^{-1}(A_i) \in \sigma(X) \otimes \sigma(Y) \subseteq \cF \otimes \cF$, we get
%\begin{align*}
%	\int_\Omega Z \bbP &= \int_{\Omega \times \Omega} f \circ (X,Y) \, \dd \Delta_\# \bbP\\
%	&= \sum_{i = 1}^M f_i \Delta_\# \bbP((X,Y)^{-1}(A_i))\\
%	&= \sum_{i = 1}^M f_i \bbP \otimes \bbP ((X,Y)^{-1}(A_i)) \\
%	&= \int_{\Omega \times \Omega} f \circ (X,Y) \, \bbP \otimes \bbP.
%\end{align*} 
%
%For general non-negative $f$ we take a monotone sequence $(f_n)_{n \ge 1}$ of simple functions that converge point-wise to $f$ and use monotone convergence twice to obtain
%\begin{align*}
%	\int_{\Omega \times \Omega} f \circ (X,Y) \, \dd \Delta_\# \bbP
%	&= \lim_{n \to \infty} \int_{\Omega \times \Omega} f_n \circ (X,Y) \, \dd \Delta_\# \bbP\\
%	&= \lim_{n \to \infty} \int_{\Omega \times \Omega} f_n \circ (X,Y) \, \bbP \otimes \bbP\\
%	&= \int_{\Omega \times \Omega} f \circ (X,Y) \, \bbP \otimes \bbP.
%\end{align*}
%
%For general $f$ we consider the positive and negative part separately and use linearity of integration.
\item (3 pts)
\textbf{2 pts}
We first observe that
\[
	\int_B Z \, \dd \bbP = \int_\Omega f \circ [(X,Y) \circ \Delta \mathbbm{1}_B] \, \dd \bbP
	= \int_{\bbR^2} f \, \dd [(X,Y) \circ \Delta \mathbbm{1}_B]_\# \bbP.
\]
On the other hand
\[
	\int_{\Omega \times B} f \circ (X,Y) \, \dd \bbP \otimes \bbP
	= \int_{\Omega \times \Omega} f \circ (X,Y)\mathbbm{1}_{\Omega \times B} \, \dd \bbP \otimes \bbP
	= \int_{\bbR^2} f \, \dd [(X,Y)\mathbbm{1}_{\Omega \times B}]_\# \bbP \otimes \bbP.
\]

Thus, the result follows if we can show that 
\[
	[(X,Y)\mathbbm{1}_{\Omega \times B}]_\# \bbP \otimes \bbP = [(X,Y)\circ \Delta \mathbbm{1}_B]_\# \bbP 
\]
as measures on $\cB^2 = \cB \otimes \cB$. 

\textbf{1 pt}
Since both measure are finite it follows from Theorem 2.15 that is suffices to check equality on the generator set $\cB \times \cB$.

\textbf{2 pts}
To this end, let $C, D \in \cB$. Then it holds that
\begin{align*}
	[(X,Y)\circ \Delta \mathbbm{1}_B]^{-1}(C \times D) 
	&= \{\omega \in \Omega \, : \, \omega \in B, X(\omega) \in C, Y(\omega) \in D\}\\
	&= X^{-1}(C) \cap Y^{-1}(D) \cap B.
\end{align*}
Recall that $B \in \sigma(Y)$ so that also $Y^{-1}(D) \cap B \in \sigma(Y)$. Thus, since $X$ and $Y$ are independent
\begin{align*}
	[(X,Y)\circ \Delta \mathbbm{1}_B]_\# \bbP(C \times D) 
	&= \bbP(X^{-1}(C) \cap Y^{-1}(D) \cap B) \\
	&= \bbP(X^{-1}(C)) \bbP(Y^{-1}(D) \cap B).
\end{align*}

\textbf{2 pts}
On the other hand, since
\[
	[(X,Y)\mathbbm{1}_{\Omega \times B}]^{-1}(C \times D) = X^{-1}(C) \times Y^{-1}(D) \cap B,
\]
we have that
\begin{align*}
	[(X,Y)\mathbbm{1}_{\Omega \times B}]_\# \bbP \otimes \bbP(C \times D)
	&= \bbP \otimes \bbP(X^{-1}(C) \times Y^{-1}(D) \cap B) \\
	&= \bbP(X^{-1}(C)) \bbP(Y^{-1}(D) \cap B). 
\end{align*}

We therefore conclude that both measures are equal on the generator set of $\cB^2$ and hence are equal on the entire \sigalg/.
\item (3 pts)
\textbf{1 pt} As noted in the proof of Lemma 11.8 the measures $(X, Y)_\# \bbP \otimes \bbP$ and $X_\# \bbP \otimes Y_\# \bbP$ agree on the generator set of $\cB^2$. Thus, using Theorem 11.5 the are equal on the entire \sigalg/. 

\textbf{2 pts} Using Fubini-Tonelli twice we then have that
\begin{align*}
	\int_{R \times C} f \dd (X, Y)_\# \bbP \otimes \bbP 
	&= \int_{R \times C} f \, \dd X_\# \bbP \otimes Y_\# \bbP\\
	&= \int_\bbR \int_C f(x,y) Y_\# \bbP(\dd y)  X_\#\bbP(\dd x) \\
	&= \int_C \int_\bbR f(x,y) X_\# \bbP(\dd X) Y_\# \bbP(\dd y).
\end{align*}

\end{enumerate}


